% !TEX root = ../main-es.tex
% Las cajas de ayuda no son contenido: texcount las ignora (ver `make wordcount`).
%TC:macro \limite [ignore,ignore]
%TC:macro \nota [ignore]
%TC:macro \TODO [ignore]
%TC:macro \parencite [ignore]
%TC:macro \textcite [ignore]
\section{Aportes de las referencias principales}
\limite{0}{Aportes de las referencias principales: seleccione dos referencias clave y
describa en qué medida contribuyen a la construcción del proyecto (por ejemplo:
fundamentos teóricos, antecedentes, metodologías, aplicaciones).}

\subsection*{\textcite{russellNorvig2020} — fundamentos teóricos}

Es el texto de la asignatura y, a la vez, el andamiaje conceptual completo del
proyecto. De él salen las tres piezas que lo estructuran: la taxonomía de agentes y
entornos del capítulo 2, que fija las seis arquitecturas comparadas y la clasificación
del entorno de tarea (parcialmente observable, estocástico, secuencial, dinámico); la
búsqueda informada del capítulo 3, cuyo A* es el planificador que el agente invoca en
cada reescritura —contenido del curso usado como motor funcional, no como cita—; y la
planeación en línea del capítulo 11 (monitoreo de acción, de plan y de objetivo), que
define qué significa detectar que un plan dejó de ser válido, el evento que dispara
toda la dinámica de replanificación.

\subsection*{\textcite{zilberstein2011metareasoning} — el marco de la pregunta}

Si R\&N describe cómo planifica un agente, Zilberstein describe cuándo le conviene
pensar más: la tradición de metarazonamiento y racionalidad acotada, donde deliberar
es una acción con costo que compite con actuar. El proyecto es exactamente ese marco
llevado a una decisión de ingeniería: el costo $c$ de edición y el precio $\kappa$ por
tick de deliberación operacionalizan el costo de cómputo, el presupuesto $B$ del agente
gobernado es una política de deliberación explícita, y el \emph{thrashing} medido por
las ediciones desperdiciadas es el modo de fallo clásico de una política de
metarazonamiento mal calibrada. La referencia aporta, además, el vocabulario para
formular la hipótesis central: existe una frontera en el plano $(\lambda, c)$ que
separa los regímenes donde conviene comprometerse con el plan de aquellos donde
conviene reescribirlo.
